\documentclass[12pt]{article}

\usepackage[utf8]{inputenc}
\usepackage[margin=1in]{geometry}
\usepackage{setspace}
\usepackage{url}
\usepackage{mathptmx}

\doublespacing
\setlength{\emergencystretch}{3em}

\title{Family Pressure, Identity, and Belonging in Alice Munro's ``Boys and Girls'' and Madeleine Thien's ``Simple Recipes''}

\author{%
  Ma Toan Bach\\
  Student Number: 112708227\\
  LSO 100--Canadian Short Story (KCC)\\
  Comparison Essay
}

\date{July 31, 2026}

\begin{document}

\maketitle

Family is often imagined as the first place where a person learns belonging, but Canadian short fiction frequently shows that family can also be where identity becomes limited, tested, or divided. Alice Munro's ``Boys and Girls'' and Madeleine Thien's ``Simple Recipes'' both focus on narrators who remember childhood experiences with a father and use those memories to understand how family shaped them. In both stories, the father is connected to work, skill, and family identity. However, that connection becomes painful when each narrator realizes that paternal love is inseparable from authority, expectation, and control. Through retrospective narration, domestic symbolism, and contrasting portrayals of paternal authority, Munro and Thien show that family gives their narrators a sense of identity while making that belonging conditional; Munro exposes the quiet internalization of gender inequality, whereas Thien portrays the lasting emotional damage caused by cultural pressure and paternal violence.

Both stories first present the father as someone whose work gives the narrator a sense of recognition. In ``Boys and Girls,'' the narrator is proud to help her father on the fox farm. She carries water, works outside, and believes that this labour makes her useful in a world that feels more important than the kitchen. Her father's praise matters deeply because when he introduces her as his ``new hired hand,'' she turns away ``red in the face with pleasure'' (Munro, 1968, p.~5). This moment shows that the narrator's identity is built partly through being seen by her father. She does not simply want to help; she wants her work to count. Nischik's chapter frames ``Boys and Girls'' through the problem of ``(un-)doing gender,'' which is useful because the narrator's pride in outdoor work challenges the feminine role that others expect her to accept (Nischik, 2007, p.~203). In Thien's ``Simple Recipes,'' the father is also introduced through skill, but the setting is the kitchen rather than the farm. The narrator remembers that her father taught her the ``simple recipe for making rice'' and that he washed the grains with care, removing ``tiny imperfections'' (Thien, 2002, p.~1). Like Munro's narrator, Thien's narrator receives identity through shared work with her father. The difference is that Munro connects the father's world to public usefulness and masculine labour, while Thien connects the father's world to food, memory, and immigrant family tradition. In both stories, then, the skills passed from father to daughter initially create belonging, but they also make the later exercise of paternal authority more emotionally damaging.

Although both fathers first seem to offer belonging, each story shows that fatherly authority also creates limits. In Munro's story, the narrator's place beside her father is temporary because other people see her gender before they see her ability. The salesman immediately weakens the father's praise by saying, ``I thought it was only a girl'' (Munro, 1968, p.~5). Later, the narrator overhears her mother say that when Laird is older he will be ``a real help,'' while the narrator can be used more in the house (Munro, 1968, p.~6). These comments teach the narrator that her labour outside is not secure. Even though she is capable, her family imagines her future as domestic. In ``Simple Recipes,'' the father's authority becomes frightening in a more direct way. At dinner, when the brother refuses the fish and challenges the father, the father responds with violence, saying, ``I don't know what kind of son you are. To be so ungrateful'' (Thien, 2002, p.~4). By turning the son's rejection of food into an act of personal betrayal, the father makes cultural loyalty dependent upon obedience to his authority. Ty's chapter frames East and Southeast Asian Canadian literature through the tension between ``the strange and the familiar,'' a useful idea for reading Thien's dinner scene because familiar family food becomes strange when it is tied to anger, rejection, and violence (Ty, 2016, p.~564). Munro's father limits the narrator through gender expectation; Thien's father damages belonging through anger, cultural pressure, and physical force.

The narrators' struggles with identity become clearest when each realizes that family belonging may require a painful kind of obedience. In ``Boys and Girls,'' the narrator begins to understand that the word ``girl'' is not neutral. She says that it once seemed ``innocent and unburdened,'' but later becomes ``a definition, always touched with emphasis, with reproach and disappointment'' (Munro, 1968, p.~8). This sentence is central because it shows that identity is being imposed from outside. The narrator is not simply growing up; she is being trained to accept a role. Her decision to open the gate for Flora becomes a quiet act of resistance. She admits, ``I did not make any decision to do this, it was just what I did'' (Munro, 1968, p.~13). By helping Flora escape, even briefly, the narrator acts against the practical world of her father and identifies with a freedom that is being taken from her. In Thien's story, the narrator's conflict is not a single act of resistance but a painful divided loyalty. After witnessing her father's violence, she cannot simply stop loving him, but she also cannot continue loving him innocently. She says that if she turns toward him, she will be ``complicit, accepting a portion of guilt'' (Thien, 2002, p.~5). This fear shows how family belonging can become morally complicated. Munro's narrator fears being reduced to a gender role; Thien's narrator fears that love itself may make her share responsibility for violence.

The two stories also use memory and first-person narration to show that identity is understood after the fact. Neither narrator fully understands the meaning of events while they are happening. Munro's narrator describes childhood with the clarity of an older voice, so the reader can see the process by which social expectations become internal beliefs. At the end of the story, after Laird exposes her role in Flora's escape, the father dismisses her action by saying, ``She's only a girl'' (Munro, 1968, p.~16). The narrator's response is even more painful: ``Maybe it was true'' (Munro, 1968, p.~16). The irony is that the sentence sounds simple, but it records a major inward defeat: the narrator begins to repeat the very judgment that has reduced her. Thien's narrator also looks backward, but the tone is more openly traumatic. She remembers her father as a magician in the kitchen, yet she also remembers that ``this violence will turn all my love to shame and grief'' (Thien, 2002, p.~5). The retrospective voice preserves both versions of the father---the careful teacher and the violent authority figure---without allowing either memory to erase the other. In Munro, memory reveals the gradual learning of gender inequality; in Thien, memory returns to trauma in order to understand why love, shame, and fear remain connected.

The settings of the stories strengthen this comparison because both homes are divided spaces. In ``Boys and Girls,'' outside work represents freedom, importance, and closeness to the father, while the kitchen represents the domestic future the narrator wants to avoid. She describes housework as ``endless, dreary, and peculiarly depressing,'' while outdoor work in her father's service seems ``ritualistically important'' (Munro, 1968, p.~6). The physical division between outside and inside becomes a symbolic division between possible identities. In ``Simple Recipes,'' the kitchen is more complicated. It begins as a place of warmth, food, and cultural connection, but it becomes the scene of conflict and later the memory of violence. The narrator says she misses the way the family once sat together while her father unveiled ``plate after plate'' (Thien, 2002, p.~2), but that same family table becomes a battleground when the brother refuses the food and the father loses control. Thien's contrast between the careful preparation of rice and the father's uncontrolled violence makes the kitchen a symbol of both cultural continuity and family fracture.

Munro's ``Boys and Girls'' and Thien's ``Simple Recipes'' therefore develop a similar theme through different literary methods. Munro uses spatial contrast, gendered language, and an ironic final sentence to show how inequality becomes internalized. Thien uses food imagery, domestic memory, and a fractured portrait of the father to show how love can survive beside shame and fear. Both stories show that identity is shaped through ordinary family experiences: farm work, cooking, eating, obedience, and memory. By looking back on childhood, both narrators discover that identity is not formed only by private choice. It is also formed by the expectations, silences, and contradictions of the family that first teaches a person where they belong.

\clearpage

\begin{center}
\textbf{References}
\end{center}

\hangindent=0.5in\noindent Munro, A. (1968). \textit{Boys and girls} [PDF]. \url{https://jerrywbrown.com/wp-content/uploads/2014/02/Boys-and-Girls-by-Alice-Munro.pdf}

\hangindent=0.5in\noindent Nischik, R. M. (2007). (Un-)doing gender: Alice Munro, ``Boys and Girls'' (1964). In R. M. Nischik (Ed.), \textit{The Canadian short story: Interpretations} (pp. 203--218). Camden House. \url{https://doi.org/10.1515/9781571136886-016}

\hangindent=0.5in\noindent Thien, M. (2002). \textit{Simple recipes} [PDF]. MostlyFiction. \url{http://mostlyfiction.com/excerpts/simplerecipes.htm}

\hangindent=0.5in\noindent Ty, E. (2016). (East and Southeast) Asian Canadian literature: The strange and the familiar. In C. Sugars (Ed.), \textit{The Oxford handbook of Canadian literature} (pp. 564--582). Oxford University Press. \url{https://doi.org/10.1093/oxfordhb/9780199941865.013.31}

\end{document}
